\documentclass{article}
\usepackage{amsmath, amssymb, amsfonts}
\usepackage{geometry}
\usepackage{booktabs}
\usepackage{xcolor}
\geometry{margin=1in}
\title{Skybox Ray Math Audit}
\author{Rune}
\date{2026-05-11}

\begin{document}
\maketitle

\section{Executive Summary}

The skybox ray computation in \texttt{Engine.hs:computeSkyboxRays} has two bugs
that together cause incorrect sky orientation (diagonal tilt) relative to world
geometry. The root cause is that the view matrix is \textbf{transposed} before
being passed to the function, but the \texttt{worldRot} extraction pattern was
designed for the non-transposed matrix. A secondary bug inverts the horizontal
ray component.

\section{Bug 1 (Root Cause): \texttt{worldRot} extracts $R$ instead of $R^{-1}$}

\subsection{The transpose that breaks everything}

At \texttt{Engine.hs:680}:
\begin{verbatim}
view = Linear.Matrix.transpose $ Camera.unViewMatrix (Camera.toMatrix camera)
\end{verbatim}

So \texttt{computeSkyboxRays} receives $V^T$, not $V$.

\subsection{What the extraction actually produces}

The extraction at line 1387:
\begin{verbatim}
worldRot = V3 (V3 v00 v10 v20) (V3 v01 v11 v21) (V3 v02 v12 v22)
\end{verbatim}

This takes columns of the input matrix and makes them \textbf{rows} of
\texttt{worldRot}.

For the \textbf{non-transposed} $V$ (what the original analysis assumes):

Columns of $V$ are:
\begin{itemize}
\item Column 0: $(s_x,\; u'_x,\; -f_x)$
\item Column 1: $(s_y,\; u'_y,\; -f_y)$
\item Column 2: $(s_z,\; u'_z,\; -f_z)$
\end{itemize}

Placing these as rows gives:
\[
\mathbf{worldRot} = \begin{pmatrix}
s_x  & u'_x & -f_x \\
s_y  & u'_y & -f_y \\
s_z  & u'_z & -f_z
\end{pmatrix} = R^T = R^{-1} \quad \text{(view $\to$ world)} \quad \checkmark
\]

For the \textbf{transposed} $V^T$ (what the code actually receives):

Columns of $V^T$ are the rows of $V$:
\begin{itemize}
\item Column 0: $(s_x,\; s_y,\; s_z) = \mathbf{s}$
\item Column 1: $(u'_x,\; u'_y,\; u'_z) = \mathbf{u'}$
\item Column 2: $(-f_x,\; -f_y,\; -f_z) = -\mathbf{f}$
\end{itemize}

Placing these as rows gives:
\[
\mathbf{worldRot} = \begin{pmatrix}
s_x  & s_y  & s_z  \\
u'_x & u'_y & u'_z \\
-f_x & -f_y & -f_z
\end{pmatrix} = R \quad \text{(world $\to$ view)} \quad \color{red}\boldsymbol{\times}
\]

The code applies the \textbf{world$\to$view} rotation to view-space directions,
producing meaningless results. For general camera orientations, $R \ne R^{-1}$.

\subsection{Numerical verification}

Camera at $(1.4, -0.7, 1.2)$ looking at origin, up $= (0, 1, 0)$:
\begin{align*}
\mathbf{f} &\approx (-0.71,\; 0.35,\; -0.61) \\
\mathbf{s} &\approx (0.65,\; 0,\; -0.76) \\
\mathbf{u'} &\approx (0.27,\; 0.93,\; 0.23)
\end{align*}

With the current code ($R$ applied to center direction $(0, 0, -1)$):
\[
R \cdot (0, 0, -1)^T = (-s_z,\; -u'_z,\; f_z) = (0.76,\; -0.23,\; -0.61)
\]

Correct result (should be $\mathbf{f}$):
\[
(-0.71,\; 0.35,\; -0.61)
\]

The center ray is $\sim\!104°$ off from the correct forward direction.

\subsection{Fix}

Since the input is $V^T$, copy rows directly instead of transposing columns:

\begin{verbatim}
worldRot = V3 (V3 v00 v01 v02) (V3 v10 v11 v12) (V3 v20 v21 v22)
\end{verbatim}

\section{Bug 2: \texttt{ndcToDir} X component has wrong sign}

\subsection{Current code}
\begin{verbatim}
viewDir = V3 (-x / fx) (y / fy) (-1)
\end{verbatim}

\subsection{Derivation of correct signs}

The projection (with Y-flip) maps view$\to$clip:
\begin{align*}
x_\text{ndc} &= f_x \cdot x_\text{view} / (-z_\text{view}) \\
y_\text{ndc} &= f_{y,\text{abs}} \cdot y_\text{view} / z_\text{view}
\end{align*}

At $z_\text{view} = -1$ (forward in right-handed view space):
\begin{align*}
x_\text{view} &= x_\text{ndc} / f_x \\
y_\text{view} &= -y_\text{ndc} / f_{y,\text{abs}}
\end{align*}

The code reads $f_y$ from the transposed projection where $f_y = -f_{y,\text{abs}}$
(negative, Y-flip baked in). So $y / f_y = -y / f_{y,\text{abs}}$ which is
\textbf{correct} for Y. The double negative cancels.

But $f_x$ is positive (no X-flip), so the X component should be $+x / f_x$:
\begin{verbatim}
viewDir = V3 (x / fx) (y / fy) (-1)
\end{verbatim}

\subsection{Why the original analysis got it wrong}

The document claims $s$ ``actually points LEFT in world space.'' This is
incorrect: $s = \text{normalize}(\mathbf{f} \times \mathbf{u})$ gives the
\textbf{right} vector in right-handed coordinates. For the given camera,
$s \approx (0.65, 0, -0.76)$ which points rightward. The erroneous claim
about $s$'s direction was used to justify the $-x$ negation.

\section{Doc Error: View Space Handedness}

Section 2.2 of the original analysis claims linear's \texttt{lookAt} produces
a \textbf{left-handed} view space. The third row being $-\mathbf{f}$ confirms
\textbf{right-handed}: objects in front of the camera get negative $Z$
(standard OpenGL convention). The claim ``$+X$ corresponds to left in world
space'' is wrong --- $+X = \mathbf{s}$ = right.

\section{Doc Error: Verification Circular}

Section 3.4's ``screen-right $\to$ world-right'' verification relies on $s$
being left. Since $s$ is right, the conclusion is invalid. With corrected
$R^{-1}$ and $+x/f_x$, screen-right correctly maps to world-right via:
\[
\frac{1}{f_x}\mathbf{s} + \mathbf{f}
\]

\section{Complete Fix}

Two changes in \texttt{src/Graphics/Haskan/Engine.hs}:

\begin{enumerate}
\item \textbf{Line 1387} --- fix \texttt{worldRot} extraction:
\begin{verbatim}
worldRot = V3 (V3 v00 v01 v02) (V3 v10 v11 v12) (V3 v20 v21 v22)
\end{verbatim}

\item \textbf{Line 1400} --- fix X sign:
\begin{verbatim}
viewDir = V3 (x / fx) (y / fy) (-1)
\end{verbatim}
\end{enumerate}

\end{document}
